\subsection{Flujo de compresión y descompresión}

Sobre cada bloque, la compresión se concibe como un flujo de etapas reversibles. El bloque de amplitudes, representado como un arreglo contiguo en
doble precisión, puede someterse primero a una transformación previa sobre su representación binaria, por ejemplo una variante de \emph{shuffle} o
\emph{bitshuffle}, con el fin de exponer regularidades que favorezcan la compresión. A continuación, el bloque transformado se entrega al
compresor sin pérdida, que genera la representación persistente utilizada por el backend de almacenamiento.

El proceso inverso ocurre únicamente cuando una operación requiere acceder al contenido del bloque. En ese momento, el bloque comprimido se
recupera, se descomprime, se invierte la transformación previa y se reconstruye el arreglo denso de amplitudes listo para ser manipulado por la
lógica del simulador. La reversibilidad de este flujo garantiza que el bloque obtenido tras la descompresión coincida exactamente con el bloque
original antes de ser comprimido, condición indispensable para preservar igualdad exacta frente a la referencia no comprimida.
